\documentclass[12pt]{article}
\usepackage{amsmath, amssymb}
\usepackage{geometry}
\usepackage{booktabs}
\geometry{margin=1in}

\title{CSE400 -- Fundamentals of Probability in Computing}
\author{Lecture Scribe: Joint Random Variables \& Joint Distributions (L15--L18)}
\date{}

\begin{document}

\maketitle

\section{Joint Random Variables}

Often two random variables must be analyzed simultaneously because one may influence the other.

Let

\[
X = \text{Random Variable 1}
\]

\[
Y = \text{Random Variable 2}
\]

The pair

\[
(X,Y)
\]

is called a \textbf{joint random variable}. The goal is to study the probability of combinations of values taken by both variables.

Examples include:

\begin{itemize}
\item Traffic speed and traffic density
\item Signal-to-noise ratio and network throughput
\item Rainfall and temperature
\item Wind speed and drone delivery time
\item Dice experiment
\end{itemize}

These situations require \textbf{joint probability distributions}.

\section{Discrete Joint Distribution -- Joint PMF}

If $X$ and $Y$ are \textbf{discrete random variables}, their joint distribution is represented using the \textbf{Joint Probability Mass Function (Joint PMF)}.

\subsection{Definition}

\[
p_{X,Y}(x,y) = P(X=x, Y=y)
\]

This represents the probability that:

\begin{itemize}
\item $X$ takes value $x$
\item AND
\item $Y$ takes value $y$
\end{itemize}

at the same time.

\subsection{Properties of Joint PMF}

\textbf{1. Non-negativity}

\[
p_{X,Y}(x,y) \ge 0
\]

\textbf{2. Total probability equals 1}

\[
\sum_x \sum_y p_{X,Y}(x,y) = 1
\]

\textbf{3. Probability of an event}

For any set $A$,

\[
P((X,Y)\in A)=\sum_{(x,y)\in A} p_{X,Y}(x,y)
\]

\section{Joint PMF Table Representation}

For discrete variables, the joint PMF can be represented using a table.

\begin{center}
\begin{tabular}{c|ccc}
\toprule
$Y \backslash X$ & $x_1$ & $x_2$ & $x_3$ \\
\midrule
$y_1$ & $p(x_1,y_1)$ & $p(x_2,y_1)$ & $p(x_3,y_1)$ \\
$y_2$ & $p(x_1,y_2)$ & $p(x_2,y_2)$ & $p(x_3,y_2)$ \\
$y_3$ & $p(x_1,y_3)$ & $p(x_2,y_3)$ & $p(x_3,y_3)$ \\
\bottomrule
\end{tabular}
\end{center}

Each entry represents

\[
P(X=x_i, Y=y_j)
\]

The sum of all entries must equal 1.

\section{Example: Rolling Two Dice}

Experiment: roll two fair dice.

Define random variables:

\[
X = \text{value on the first die}
\]

\[
T = \text{total (sum) of both dice}
\]

Possible values:

\[
X \in \{1,2,3,4,5,6\}
\]

\[
T \in \{2,3,4,5,6,7,8,9,10,11,12\}
\]

Total sample outcomes:

\[
6 \times 6 = 36
\]

Each outcome has probability

\[
\frac{1}{36}
\]

\subsection{Joint PMF}

\[
p_{X,T}(x,t)=P(X=x,T=t)
\]

Example calculations:

\textbf{Case 1}

\[
P(X=1,T=2)
\]

Outcome: $(1,1)$

\[
P=\frac{1}{36}
\]

\textbf{Case 2}

\[
P(X=1,T=7)
\]

Outcome: $(1,6)$

\[
P=\frac{1}{36}
\]

\textbf{Case 3}

\[
P(X=3,T=7)
\]

Outcome: $(3,4)$

\[
P=\frac{1}{36}
\]

If a pair $(x,t)$ cannot occur:

\[
p_{X,T}(x,t)=0
\]

Example:

\[
P(X=6,T=3)=0
\]

\section{Continuous Case -- Joint PDF}

If $X$ and $Y$ are \textbf{continuous random variables}, we use the \textbf{Joint Probability Density Function}.

\subsection{Definition}

\[
f_{X,Y}(x,y)
\]

The probability that $(X,Y)$ lies in region $R$ is

\[
P((X,Y)\in R)=\int\int_R f_{X,Y}(x,y)\,dx\,dy
\]

\subsection{Properties}

\textbf{Non-negativity}

\[
f_{X,Y}(x,y) \ge 0
\]

\textbf{Total probability}

\[
\int_{-\infty}^{\infty}\int_{-\infty}^{\infty} f_{X,Y}(x,y)\,dx\,dy = 1
\]

\section{Geometric Interpretation}

In continuous distributions, probability corresponds to the \textbf{volume under a surface}.

Example probability:

\[
P(a<X<b, c<Y<d)
\]

\[
=\int_a^b \int_c^d f_{X,Y}(x,y)\,dy\,dx
\]

\section{Worked Example: Normalization}

Suppose

\[
f_{X,Y}(x,y)=c
\]

inside region $R$.

\textbf{Step 1}

\[
\int\int_R f_{X,Y}(x,y)\,dx\,dy =1
\]

\textbf{Step 2}

\[
\int\int_R c\,dx\,dy =1
\]

\textbf{Step 3}

\[
c \int\int_R dx\,dy =1
\]

If

\[
a<x<b
\]

\[
c<y<d
\]

Area:

\[
(b-a)(d-c)
\]

\textbf{Step 4}

\[
c(b-a)(d-c)=1
\]

\[
c=\frac{1}{(b-a)(d-c)}
\]

\section{Joint Cumulative Distribution Function}

\subsection{Definition}

\[
F_{X,Y}(x,y)=P(X\le x, Y\le y)
\]

\section{Continuous Case}

\[
F_{X,Y}(x,y)=\int_{-\infty}^{x}\int_{-\infty}^{y} f_{X,Y}(u,v)\,dv\,du
\]

\subsection{Relation Between CDF and PDF}

\[
f_{X,Y}(x,y)=
\frac{\partial^2}{\partial x \partial y}
F_{X,Y}(x,y)
\]

\section{Discrete Case}

\[
F_{X,Y}(x,y)=\sum_{i\le x}\sum_{j\le y} p_{X,Y}(i,j)
\]

\section{Properties of Joint CDF}

\textbf{Range}

\[
0 \le F_{X,Y}(x,y) \le 1
\]

\textbf{Limits}

\[
\lim_{x\to -\infty}F(x,y)=0
\]

\[
\lim_{y\to -\infty}F(x,y)=0
\]

\textbf{Upper limit}

\[
\lim_{x,y\to\infty}F(x,y)=1
\]

\textbf{Monotonicity}

The CDF is non-decreasing in both variables.

\subsection{Rectangle Probability Formula}

\[
P(a<X\le b, c<Y\le d)
\]

\[
=F(b,d)-F(a,d)-F(b,c)+F(a,c)
\]

\section{Key Formulas for Exam Revision}

\subsection{Joint PMF}

\[
p_{X,Y}(x,y)=P(X=x,Y=y)
\]

\[
\sum_x\sum_y p_{X,Y}(x,y)=1
\]

\subsection{Joint PDF}

\[
P((X,Y)\in R)=\int\int_R f_{X,Y}(x,y)\,dx\,dy
\]

\[
\int\int f_{X,Y}(x,y)\,dx\,dy=1
\]

\subsection{Joint CDF}

\[
F_{X,Y}(x,y)=P(X\le x,Y\le y)
\]

Continuous case:

\[
F(x,y)=\int_{-\infty}^{x}\int_{-\infty}^{y}f(u,v)\,dv\,du
\]

PDF from CDF:

\[
f(x,y)=\frac{\partial^2}{\partial x \partial y}F(x,y)
\]

\end{document}
